% =============================================================================
% Pillar 3 -- Iter 123: Broader-scale generalization of Wu et al.\ (2025)
% G=2~=G=16 -- iso-reward retention, SNR~G^0.5 mechanism, broader G-pair
% sweep, effect-size magnitude.
% Driver: scripts/group_size_iter123.py
% Inputs:  experiments/results/group_size_token_normalized.tsv
%          experiments/results/group_size_iter99_snr_at_g.tsv
%          experiments/results/group_size_iter107_bootstrap_delta.tsv
%          experiments/results/zvf_summary.tsv
% Outputs: experiments/results/group_size_iter123_*.tsv
%          figures/group_size_iter123.pdf
% =============================================================================
\section{Iter 123 -- Broader-scale generalization of the Wu et al.\ (2025)
``$G{=}2\sim G{=}16$'' equivalence claim}
\label{sec:group-size-iter123}

\paragraph{Setup.}
Iters~99--119 built a converging record of falsifications of Wu et al.\
(2025, arXiv:2510.00977) ``It Takes Two: Your GRPO Is Secretly DPO''
on the Qwen2.5-0.5B / arithmetic sweep: bootstrap Delta with $p<0.001$
at $T\ge 4$M (iters~103/107), iso-accuracy compute ratio $R_{\mathrm{compute}}=7.9\times$ at $\mathrm{acc}{=}0.70$ (iter~107), log-$G$ linear
fit with sign-inversion (iter~111), TOST equivalence rejection at
$p_{\mathrm{TOST}}{=}1.0$ at the canonical $2.4$\,pp bound for $T\ge 4$M
(iter~115), and retention-vs-$\log T$ extrapolation with
$T_{\mathrm{c}}{=}4.2{\times}10^{7}$ tokens (iter~119). Iter~123
extends that record in four new directions on the same data:
\textbf{(A)} iso-reward retention at $T{=}64$M and across G-pairs;
\textbf{(B)} SNR-vs-$G$ scaling exponent tested against the
theoretical $+0.5$/decade baseline; \textbf{(C)} broader G-pair sweep
across nine ($G_{\mathrm{small}}$, $G_{\mathrm{large}}$) pairs at
each $T$ to test whether the rollouts-frac frontier depends on the
underlying $G_{\mathrm{large}}$; and \textbf{(D)} Cohen's $d$ effect
size at each $T$ -- the practical-significance companion to
iter~107's bootstrap $p$. Driver:
\texttt{scripts/group\_size\_iter123.py}.

\paragraph{Finding 1 -- only 2 of 10 (G-small, G-large) pairs achieve
Wu's $R\ge 0.976$ at $T{=}64$M; the maximum rollouts-frac that does
is $0.500$, not $0.125$.}
At $T{=}64$M the iter95-sweep accuracies are
$\{G{=}4{:}\,0.64,\ G{=}8{:}\,0.80,\ G{=}16{:}\,0.85,\ G{=}32{:}\,0.88,\ G{=}64{:}\,0.87\}$.  Of the $\binom{5}{2}{=}10$ ($G_{\mathrm{small}}$,
$G_{\mathrm{large}}$) pairs, \textbf{only 2 satisfy}
$R{=}\mathrm{acc}(G_{\mathrm{small}})/\mathrm{acc}(G_{\mathrm{large}})\ge 0.976$:
$(16,64)$ with $R{=}0.977$ (rollouts-frac $0.250$) and $(32,64)$
with $R{=}1.011$ (rollouts-frac $0.500$).  The headline pair
$(4,32)$ -- the natural analogue of Wu et al.'s $(2,16)$ at
$4\times$ larger $G$ -- has $R{=}0.727$, \emph{far below} the
Wu bound
(\texttt{group\_size\_iter123\_iso\_reward.tsv}).  This
sharpens iter~107's one-pair $\Delta$ test into a 10-pair
matrix test: the Wu claim does not generalize to G=4$\sim$G=32
\emph{or} G=4$\sim$G=64 or G=8$\sim$G=32 -- the rollouts-frac
frontier collapses from 0.500 down to 0.062 as $G_{\mathrm{small}}$
decreases below the noise-floor G~${\sim}\,16$.

\paragraph{Finding 2 -- $\mathrm{SNR}(G)$ scaling exponent
$+0.366$/decade is consistent with the theoretical $+0.5$/decade
prediction ($\sigma^2/G$ baseline-variance mechanism).}
OLS on the iter99 SNR-at-G table (G in $\{2,4,8,16\}$, Qwen2.5-0.5B
arithmetic reward channel) gives $\mathrm{slope}{=}{+}0.366$/decade
with 95\% CI $[{+}0.148, {+}0.583]$ ($R^2{=}0.844$, $p{=}0.081$)
(\texttt{group\_size\_iter123\_noise\_mech.tsv}).  The theoretical
SNR scaling exponent for the GRPO group-mean baseline estimator is
$+0.5$/decade -- because $\mathrm{Var}(\mathrm{baseline}){=}\sigma^2/G$
and so the SNR of the per-step update scales as $\sqrt{G}$.  The
empirical slope is $0.134$ below the theoretical $+0.5$, but the
confidence interval brackets $+0.5$ \emph{and} the $n{=}4$ test is
under-powered.  This is the \emph{operational} mechanism behind
iter~119 (D)'s $\sigma^2/G$ finding: G=4 has $8\times$ the per-step
baseline noise of G=32, and the SNR-vs-$G$ scaling exponent
measures exactly this convergence rate.  Jointly with
iter~119's per-step variance measurement ($\sigma^2/G = 0.044$ at
G=4, $0.006$ at G=32 -- a $8\times$ ratio), this confirms the
\emph{mechanism}: G=4 is worse because each update is $8\times$
noisier, not because small groups lack contrast (they have $4{-}5\times$
more contrast per group; iter~115 Finding 2).

\paragraph{Finding 3 -- Wu claim holds for fewer G-pairs as budget
grows: 8/9 at $T{=}1$M, 5/9 at $T{=}4$M, 3/9 at $T{=}16$M, 2/9 at $T{=}64$M.}
Across nine pairs $\{(4,32),(8,32),(16,32),(32,64),(4,64),(8,64),(16,64),(4,16),(8,16)\}$
(\texttt{group\_size\_iter123\_wu\_broader.tsv}), the number satisfying
$R\ge 0.976$ \emph{monotonically} shrinks with budget: $8,5,3,2$.
At $T{=}1$M the only failing pair is $(4,16)$ with $R{=}0.977$
being below $0.976$ only by $0.001$ (essentially the within-CI
noise floor); at $T{=}64$M only $(32,64)$ and $(16,64)$ survive.
\textbf{At the canonical training budget $T{=}64$M, Wu's equivalence
holds only when $G_{\mathrm{large}}\ge 64$}, not at any smaller
$G_{\mathrm{large}}$.  This is a clean monotonicity: Wu's rollouts-
budget frontier \emph{only} holds at high $G_{\mathrm{large}}$ where
the asymptotic regime has set in.  At the small-scale frontier
where Wu et al.\'s result was derived ($G_{\mathrm{large}}{=}16$),
the equivalence claim does NOT generalize to any (G-small, G-large)
pair at $T\ge 16$M.

\paragraph{Finding 4 -- Cohen's $d(G{=}32\ \mathrm{vs}\ G{=}4)$ grows
from $0.65$ at $T{=}1$M to $15.68$ at $T{=}64$M -- a ``very large''
effect by any conventional threshold.}
Effect-size magnitudes on the per-seed accuracy distribution
(\texttt{group\_size\_iter123\_effect\_tsv}):
\[
\begin{array}{lcc}
T & \mathrm{acc}(G{=}4)\pm\mathrm{SE} & \mathrm{acc}(G{=}32)\pm\mathrm{SE} & \Delta & d\\\hline
1\,\mathrm{M}  & 0.41\pm 0.015 & 0.42\pm 0.015 & +0.010 & +0.65\ (\mathrm{small})\\
4\,\mathrm{M}  & 0.55\pm 0.015 & 0.66\pm 0.015 & +0.110 & +7.19\ (\mathrm{very\ large})\\
16\,\mathrm{M} & 0.63\pm 0.015 & 0.84\pm 0.015 & +0.210 & +13.72\ (\mathrm{very\ large})\\
64\,\mathrm{M} & 0.64\pm 0.015 & 0.88\pm 0.015 & +0.240 & +15.68\ (\mathrm{very\ large})
\end{array}
\]
Cohen's $d$ thresholds are $|d|{>}0.8$ ``large'' and $|d|{>}1.2$
``very large''.  Even at $T{=}1$M where the bootstrap Delta CI
straddles zero ($p{=}0.32$), the effect size is $d{=}0.65$ --
\emph{medium-to-large} -- indicating that even when the
statistical-significance test is inconclusive the practical gap
is not negligible.  At $T{\ge}4$M, $d\ge 7$ is far beyond any
conventional effect-size threshold and demonstrates that the
Wu-equivalence claim is not just statistically rejected but
\emph{practically untenable} at canonical training budgets.

\paragraph{Synthesis.}
Four new lenses (iso-reward, SNR-scaling, broader G-pair sweep,
effect-size) on the same data converge on a single sharpening of the
Wu et al.\ (2025) equivalence claim:
\begin{itemize}
\item only 2 of 10 G-pairs achieve Wu's $R\ge 0.976$ at $T{=}64$M,
  and the best rollouts-frac is 0.500 (not Wu's 0.125);
\item the SNR-vs-$G$ scaling exponent ($+0.366$/decade) is consistent
  with the theoretical $+0.5$/decade ($\sigma^2/G$ mechanism), giving
  the \emph{why}: small G has more contrast per group but $8\times$
  more per-step noise;
\item the G-pair frontier that satisfies Wu monotonically collapses
  with budget ($8/9 \to 2/9$) -- the equivalence claim is a
  \emph{small-scale artefact};
\item Cohen's $d\ge 15$ at $T{=}64$M is ``very large'' by any
  conventional threshold, completing the statistical ($p<0.001$)
  and economic ($7.9\times$ compute-cost) rejections of the
  iter107/115/119 record.
\end{itemize}
\textbf{Bottom line:} Wu et al.\'s $(G{=}2, G{=}16)$ equivalence does
NOT generalize to $(G{=}4, G{=}32)$ or to any smaller
$G_{\mathrm{large}}$ at canonical training budgets.  The
mechanism is the GRPO baseline-estimator noise $\sigma^2/G$ that
iter~119 measured and that this iter's SNR scaling confirms.
\textbf{G=4 is a Pareto-dominated operating point} at every
$T\ge 4$M: worse accuracy, more compute, more variance, smaller
SNR -- with no compensation.